\uuid{1rlX}
\titre{Raisonnements}
\auteur{bodin ; croizat ; sacre}
\organisation{exo7}
\licence{CC BY-NC-SA 4.0}
\datecreate{2022-03-08}
\niveau{L1}
\module{Algèbre}
\chapitre{Logique, ensemble, raisonnement}
\sousChapitre{Absurde et contraposée}
\difficulte{3}
\theme{Raisonnements}
\exo7qcmid{hilisit.logique.054}

\contenu{
\begin{question}
\qclassification{100.03}
Pour montrer que $4^n > 20n$ pour des entiers $n$ naturles suffisamment grands je fais une preuve par récurrence.
Je peux commencer l'initialisation avec :
\begin{answers}
	\bad{$n=0$}
	\bad{$n=1$}
	\bad{$n=2$}
	\good{$n=3$}
\end{answers}
\begin{explanations}
La propriété $4^n > 20n$ est vraie pour $n=0$ mais fausse pour $n=1$ et pour $n=2$.
Ensuite pour tout $n\ge3$ elle est vraie.
Je peux donc choisir l'initialisation avec $n=3$.\end{explanations}
\end{question}
}
